
This paper evaluated a post-hoc, generation-free hallucination detection approach applying frozen DeBERTa NLI contradiction scoring to HaluEval across 60,000 pairs. AUROC = 0.709 on dialogue and 0.644 on QA indicate detectable signal for commission-type hallucinations; AUROC = 0.530 on summarization is near chance and near the structural ceiling for contradiction-based detection on omission-type hallucinations. Wilcoxon rank-sum tests confirm score ordering by hallucination label on all three tasks.

The commission/omission boundary offers a structural account of when NLI-based post-hoc detection is viable. Three directions follow from the results: (1) omission-aware detection via coverage-based metrics for summarization tasks; (2) design ablation experiments to attribute performance to specific configuration choices (h-m2 through h-m4); and (3) cross-model generalization to determine whether the commission/omission boundary is model-invariant or specific to DeBERTa-v3-large.
